% Course notes for the LaTeX loader (W5c3, ADR-0035). The owner's
% checklist: sections by h, the table by t, a formula read as math.
\documentclass[11pt]{article}
\usepackage[english]{babel}
\usepackage{amsmath,amsthm,booktabs,hyperref}
\newcommand{\course}{Biology 101}
\newcommand{\R}{\mathbb{R}}
\newtheorem{definition}{Definition}
\title{Cells and Growth}
\author{Ada Example}
\date{Fall term}

\begin{document}
\maketitle

\begin{abstract}
Notes for \course{} on how cells divide, and the arithmetic of growth.
\end{abstract}

\section{Cells}\label{sec:cells}
Every living thing is made of cells.\footnote{Viruses are the usual exception.}
A cell has three parts we care about:
\begin{itemize}
  \item the \emph{membrane}, which holds it together;
  \item the \textbf{nucleus}, which holds the DNA;
  \item the cytoplasm, everything else.
\end{itemize}

\begin{definition}[Mitosis]\label{def:mitosis}
Mitosis is the division of one cell into two identical cells.
\end{definition}

\subsection{How fast they divide}
A population that doubles every hour grows as
\begin{equation}\label{eq:growth}
  N(t) = N_0 \cdot 2^{t}
\end{equation}
where $t$ is the time in hours and $N_0 \in \R$ is the starting count.
Equation~\eqref{eq:growth} is used again in Section~\ref{sec:data}.

\section{Data}\label{sec:data}
Table~\ref{tab:counts} gives the counts from the lab
\cite[p.~12]{doe2020}.

\begin{table}[h]
  \centering
  \caption{Cell counts by hour}\label{tab:counts}
  \begin{tabular}{lr}
    \toprule
    Hour & Cells \\
    \midrule
    0 & 100 \\
    1 & 200 \\
    2 & 400 \\
    \bottomrule
  \end{tabular}
\end{table}

% Read by notes.tex through \input: an included file in the same folder.
\section{Summary}
\begin{enumerate}
  \item Cells divide by mitosis (Definition~\ref{def:mitosis}).
  \item Growth follows $N(t) = N_0 \cdot 2^{t}$.
\end{enumerate}


\end{document}
